\section{Toward Bethlehem}

On December 24 the Vigil Mass takes us into Joseph's trouble. Mary is with
child, and Joseph is considering what to do. The angel tells him to take Mary
as his wife and to give her son the name Jesus, because he will save his
people from their sins (Matt. 1:18--21). The Collect asks that we who now
receive him gladly as Redeemer may look on him with confidence when he comes
as Judge. On the day we have almost reached the manger, the Church still asks
for the grace to look up when the Son of Man comes.

At the first Mass of Christmas, the Mass in the Night, the Roman station is at
the Crib in St Mary Major. The Gospel tells of shepherds keeping watch and of
a sign they can find for themselves. “You shall find the infant wrapped in
swaddling clothes and laid in a manger” (Luke 2:12). Mary has wrapped him and
put him where the animals feed. It is this detail, after the announcement that
the child is Christ the Lord, that the angel tells the shepherds to look for.

Gregory the Great hears in the name Bethlehem the Hebrew for “house of bread”
(Homily 8 on the Gospels, §1). Christ says that he is the living bread come
down from heaven and that the bread he will give is his flesh for the life of
the world (John 6:51--52). Gregory makes a further connection. Isaiah calls
all flesh grass, which withers, and Jesus speaks of a grain of wheat that must
fall into the earth and die to bear fruit (Isa. 40:6--8, John 12:24--25). The
Son took flesh that could suffer and die. Gregory says that in taking it he
turned our grass into wheat. At the manger the faithful find food in the flesh
of the very one who will die like the grain in the earth.

Luke tells how, at the other end of Jesus' life, Joseph of Arimathea takes his
body from the cross, wraps it in linen and lays it in a tomb (Luke 23:50--53).
At the arrest, the soldiers had bound Jesus after he went out to meet them
(John 18:4, 12). At the birth and at the burial the wrappings are acts of
care. The one whom Mary tended let himself be taken, and went to his death in
the flesh she had borne.

St Gregory of Nazianzus recalls the infant clothes when he speaks of the
Resurrection. “He was wrapped in swaddling clothes—but He took off the
swathing bands of the grave by His rising again” (Oration 29.19). The body
Mary tended is the body he raised. Gregory goes on to speak of the one laid in
a manger whom angels glorified and the Magi worshipped. The Son could truly
die in the nature he took from us, and death could not keep him.

St Cyril of Alexandria sees a table in the manger (Commentary on Luke, Sermon
1, on 2:7). He says that the Son found us living at the level of beasts. In
that very place where beasts are fed, he offers something that can restore us
to human life with God. Cyril calls the manger “even His own table,” where we
find “the bread from heaven, which is the body of life.” In the Eucharist the
flesh that gives us life is the flesh the Word made his own. The risen Christ
gives us the body he offered for us. This is what it means to come hungry to
the altar at Christmas. We need from him the life he came into our world to
give.

St Paul says that God sent his Son, born of a woman and born under the law, so
that we might receive adoption. God sends the Spirit of his Son into our
hearts, and we can call him Father (Gal. 4:4--7). St Leo takes this as a
reason to rejoice on the feast. “The birthday of the Head is the birthday of
the body” (Nativity Sermon 26.2). Each of us comes to the font at a particular
time. In baptism the whole Church is born with Christ, crucified with him and
raised with him. He has come down into our life so that we may share his. The
peace the angels declare also belongs to that new life. Leo says that Christ
brings Jew and Gentile to his Father in one Spirit (Nativity Sermon 26.5, Eph.
2:14--18). The Father whom we call upon in Christ gives us brothers and
sisters to live with in peace.

At the Mass in the Night, before the Gospel about the shepherds, St Paul tells
Titus that the grace of God has appeared. It teaches us how to live now, while
we await the appearing in glory of our Saviour Jesus Christ (Titus 2:11--15).
St Cyril of Jerusalem uses this passage when he teaches those preparing for
baptism to look for Christ's return. He asks them to greet his second coming
with the words the crowd spoke when Jesus entered Jerusalem (Catechetical
Lecture 15.1--3). We already pray them at Mass, at the end of the Sanctus. Say
them at the crib. You can say them again in thanksgiving when you have
received him at the altar. “Blessed is he that cometh in the name of the Lord”
(Matt. 21:9).
